\chapter{関連研究}

本章では，較正の評価指標，確信度を言語で表明させる手法，多言語環境での較正，較正の改善手法の順に既存研究を概観し，本研究の位置づけを示す．

\section{較正の評価指標}

較正は，モデルが出力する確信度と実際の正答率の一致を指す概念であり，もとは気象予報の評価から発展した．機械学習の分野では Guo ら \cite{guo2017calibration} が広く知られる．Guo らは，分類精度の高い深層ニューラルネットワークが，層の浅い従来のモデルに比べてむしろ較正が悪化していることを示した．すなわち精度と較正は別の性質であり，精度が高いモデルが自動的に信頼できる確信度を出すわけではない．

較正の程度を測る指標として広く用いられるのが ECE（Expected Calibration Error）であり，区間分割を用いる推定法は Naeini ら \cite{naeini2015binning} に基づく．ECE は解釈しやすい一方で，区間の数と確信度の分布に値が依存する．確信度が狭い範囲に偏っていると，標本の大半が少数の区間に入るため，区間ごとの差を平均するという計算が，全体の平均どうしの差に近づく．この状況では ECE の値だけから較正の良否と正答率の高低を読み分けにくい．本研究の予備実験でもこの状況が生じた（5.4 節）．

ECE を補う指標として，確信度と正誤の二乗誤差の平均である Brier Score と，確信度の順序のみを評価する AUROC が併用される．Murphy \cite{murphy1973partition} は Brier Score を信頼度（reliability），識別度（resolution），不確実度（uncertainty）の 3 成分に分解できることを示した．ここで注意を要するのは，不確実度が正答率のみで決まるからといって，信頼度が正答率と独立になるわけではないことである．信頼度の定義には各区間の正答率が含まれ，区間が 1 つしかない極端な場合，信頼度は平均確信度と平均正答率の差の二乗に一致する．したがって分解は，較正の水準と正答率の高低を機械的に切り分ける手段ではない．本研究では，分解を交絡を解消する手段としてではなく，Brier Score の内訳を把握するための補助的な情報として用いる．各指標の定義は 3.4 節で述べる．

\section{確信度の言語による表明}

LLM の確信度を得る素朴な方法は，モデルが割り当てる確率を用いることである．Tian ら \cite{tian2023justask} は，人間のフィードバックによる強化学習で調整されたモデルでは内部の確率に基づく較正が崩れていることを指摘した．また商用の API では対数確率が取得できない場合も多く，実用上の制約もある．

これに対し，モデル自身に確信度を言語で答えさせる手法が提案されている．Lin ら \cite{lin2022teaching} は，モデルが不確実性を言葉で表現するよう学習させられることを示し，Kadavath ら \cite{kadavath2022know} は，モデルが自身の回答の正しさをある程度予測できることを報告した．Tian ら \cite{tian2023justask} は，プロンプトのみで確信度を引き出す複数の方式を比較し，言語で表明させた確信度がサンプリングによって推定した条件付き確率より良好な較正を示す場合があることを報告した．本研究が比較する 3 方式は，この研究の枠組みを踏まえたものである．Xiong ら \cite{xiong2024llmuncertainty} は確信度の引き出し方を体系的に整理し，プロンプトの与え方によって結果が大きく変わること，そして LLM が一貫して自信過剰である傾向を示した．また Dai ら \cite{dai2026rescaling} は，0 から 100 の尺度では確信度が切りのよい数値に集中することを指摘し，より粗い尺度のほうが良好な結果を示す場合があると報告している．本研究が言語表現による方式を含めるのは，利用者に提示する際には数値より「かなり自信がある」という表現のほうが直感的に伝わる可能性がある一方，段階数が限られるため確信度の解像度が落ちる恐れもあり，この得失を実測で確認する必要があるためである．

自信過剰の原因について，Seo ら \cite{seo2025advice} は，確信度の推定が回答の内容に依存していないことが主要因であると指摘した．すなわちモデルは，自分が何と答えたかを踏まえずに確信度を述べている場合がある．この知見は，回答を得たうえで改めて確信度を尋ねる 2 段階方式を設計する動機となる．ただし，Seo らの手法自体はモデルの重みを更新するものであり，本研究が引用するのは自信過剰の原因に関する診断的な知見に限られる．回答を提示して尋ねるだけで較正が改善するとまでは言えないため，本研究はこれを検証の対象として設定する．

\section{多言語環境における較正}

多言語での較正を扱った研究として，Xue ら \cite{xue2025mlingconf} の MlingConf がある．同研究は英語，日本語，中国語，フランス語，タイ語の 5 言語を対象に，モデル内部の確率を用いる手法，p(True) を用いる手法，verbalized confidence の 3 手法を比較した．同研究の報告によれば，言語によって較正の質は異なり，英語での較正が他言語より良好な傾向がある．たとえば GPT-3.5 を用いた TriviaQA では，英語のECE が 16.52 であるのに対し日本語では 34.49 と報告されている．

MlingConf は日本語を対象に含む数少ない研究であり，本研究にとって最も近い先行研究である．本研究との違いは，比較している条件の組み合わせにある．MlingConf が扱うのは主として推定手法の違いであるのに対し，本研究は verbalized confidence の内部で，回答と確信度を同時に出すか別のターンに分けるか，および確信度を数値で答えさせるか言語表現で答えさせるかという組み合わせを比較する．

\section{較正の改善手法と本研究の位置づけ}

較正を改善する手法は，学習を伴うもの，出力を後処理するもの，プロンプトによるものに大別できる．学習を伴う手法として，Li ら \cite{li2025conftuner} の ConfTuner は確信度を適切に表明するようモデルを訓練する手法であり，Xiao ら \cite{xiao2025restoringcalibration} は事後学習によって損なわれた較正を回復させる手法を報告している．後処理による手法では温度スケーリングが代表例であり，Guo ら \cite{guo2017calibration} が分類モデルに対して有効性を示し，LLM に対しても Xie ら \cite{xie2024adaptivetemp} が適応的な温度調整を報告している．前者はモデルの重みへの接触を要し，後者は補正のパラメータを決めるための検証データを必要とするため，いずれも商用 API 経由でのみ利用できるモデルには適用しにくい．

本研究が対象とするのは，モデルの重みにも出力の後処理にも触れず，入力するプロンプトのみで較正を改善する方向である．この方向には，重みにアクセスできない商用モデルにも適用できること，追加の学習コストがかからないこと，モデルが更新されても手法自体は使い続けられることという利点がある．黒箱環境での較正手法については Xie ら \cite{xie2024blackboxsurvey} に整理がある．一方で，プロンプトによる手法は文面の細部に結果が左右されやすく，効果の大きさもモデルやタスクに依存するため，どのような設計が有効かは対象とする条件のもとで実測して確かめるほかない．

以上を踏まえ，本研究の位置づけを整理する．方式の比較そのものは Tian ら \cite{tian2023justask} やXiong ら \cite{xiong2024llmuncertainty} が英語で行っており，MlingConf は日本語を対象に含んでいる．本研究はこれらに対し，claude-sonnet-4-6 と日本語版 MGSM という具体的な条件のもとで，出力のタイミングと確信度の表現形式を組み合わせた 3 条件を同一の評価手続きで比較する．したがって本研究は，日本語での較正研究が存在しないという主張をするものではない．また 2.1 節で述べたとおり ECE の値は確信度の分布に左右され，Murphy 分解も較正と正答率を機械的に切り分けるわけではないため，複数の指標と確信度の分布を併記し，各指標から何が言えて何が言えないかを明示する．そのうえで，較正の改善をプロンプト設計の問題として扱う．この方向自体は先行研究にもあるため，新規性は改善手法の具体的な設計と，その日本語タスクでの検証に置く．
